% Figure: double-pendulum geometry and its two generalised coordinates.
% Two massless links, two bobs, angles measured from the downward
% vertical at each pivot.
\begin{figure}[htbp]
\centering
\begin{tikzpicture}[
  link/.style={draw=engblack, very thick},
  bob/.style={draw=engblue, fill=engblue, circle, inner sep=2.5pt},
  vert/.style={draw=enggray, dashed, thin},
  ang/.style={draw=engred, -{Stealth}, thick},
  label/.style={font=\small},
]
% Pivot at origin (top)
\fill[engblack] (0,0) circle (1.6pt);
\node[font=\scriptsize, anchor=south] at (0,0.1) {pivot};

% Upper link to bob 1: theta1 = 35 deg from vertical down
% bob1 at (L1 sin t1, -L1 cos t1), L1 = 3
\coordinate (b1) at ({3*sin(35)}, {-3*cos(35)});
% downward vertical guide from pivot
\draw[vert] (0,0) -- (0,-3.4);
\draw[link] (0,0) -- (b1);
\node[bob] (B1) at (b1) {};
\node[label, anchor=west] at (b1) {$m_1$};
% angle arc theta1
\draw[ang] (0,-1.0) arc (270:{270+35}:1.0);
\node[engred, font=\scriptsize] at ({0.75*sin(17.5)}, {-1.25*cos(17.5)}) {$\theta_1$};
% upper link length label
\node[font=\scriptsize, enggray, anchor=east] at ({1.5*sin(35)-0.15}, {-1.5*cos(35)}) {$L_1$};

% Lower link from bob1: theta2 = -20 deg from vertical down at b1, L2 = 2.6
\coordinate (b2) at ({3*sin(35) + 2.6*sin(-20)}, {-3*cos(35) - 2.6*cos(-20)});
\draw[vert] (b1) -- ($(b1)+(0,-3.0)$);
\draw[link] (b1) -- (b2);
\node[bob] (B2) at (b2) {};
\node[label, anchor=west] at (b2) {$m_2$};
% angle arc theta2 (negative, to the left of vertical)
\draw[ang] ($(b1)+(0,-0.9)$) arc (270:{270-20}:0.9);
\node[engred, font=\scriptsize] at ($(b1)+({-0.7*sin(10)},{-1.15})$) {$\theta_2$};
\node[font=\scriptsize, enggray, anchor=west] at ($(b1)!0.5!(b2)+(0.12,0)$) {$L_2$};

\end{tikzpicture}
\caption[Double-pendulum geometry and generalised coordinates]{The
planar double pendulum and its two generalised coordinates. Each swing
angle $\theta_1$ and $\theta_2$ is measured from the downward vertical
(grey dashed) at its own pivot: $\theta_1$ at the fixed pivot, $\theta_2$
at the first bob. The four Cartesian coordinates of the two bobs are
reduced to two by the inextensible-link constraints, which the angle
coordinates satisfy identically, so the Lagrangian carries no
constraint forces. The two coordinates make the system minimal: any
fewer cannot locate it, any more would be redundant. The chapter's
project derives the equations of motion in these coordinates two ways
and reconciles them.}
\label{fig:vol03ch05:double-pendulum}
\end{figure}
